\begin{frame}{7.3 — Definir casos antes de executar}
\small
\begin{tabularx}{\textwidth}{lXrr}
\toprule
Caso & Leituras & Limite & Esperado \\
\midrule
Normal misto & \texttt{\{28, 30, 32\}} & 30 & 2 \\
Nenhuma ocorrência & \texttt{\{25, 28\}} & 30 & 0 \\
Todas as ocorrências & \texttt{\{30, 31\}} & 30 & 2 \\
Igualdade & \texttt{\{30\}} & 30 & 1 \\
Vazio & Nenhuma leitura & 30 & 0 \\
Valores negativos & \texttt{\{-2, 0, 2\}} & 0 & 2 \\
\bottomrule
\end{tabularx}
\vspace{0.3cm}
O esperado deve vir da especificação, não da execução do próprio código que será testado.
\end{frame}
